\chapter{What Staff means}
\label{ch:staff}

\section*{What this chapter is for}

Staff is not Senior with more years on it, and a candidate who prepares as a
very good Senior will be assessed as a very good Senior. The difference is
legible in published ladders, it is specific, and it changes what you should
choose to talk about.

\section{One word does most of the work}

Published engineering ladders are worth reading precisely because they are
published: the wording can be quoted rather than paraphrased. The most cited of
them describes Staff engineers as taking the lead on directing and implementing
solutions to significantly complex, \reported{\textbf{unscoped}
problems}{etsy-ladder}, whose work extends across many parts of a domain and
may begin to influence a whole department.

\emph{Unscoped} is the load-bearing word.

\begin{kitnote}
A Senior engineer is given a problem and is excellent at solving it. A Staff
engineer decides which problem is worth solving, and is then accountable for
having chosen it.

Everything else in this chapter follows from that sentence. The scope question,
the influence question, the ambiguity question and the project deep-dive are
four ways of asking it.
\end{kitnote}

Read more than one ladder before you use any of them; several companies publish
theirs and an aggregator collects them.\source{progression-fyi} They disagree in
detail and agree on this.

\section{Scope, not effort}

The most common way to be doing Staff-shaped work and not get credit for it is
to be doing the coordination that makes a team function \dash{} the following-up,
the documentation, the unblocking, the noticing \dash{} and to have none of it
written into any ladder. That failure mode has a canonical
treatment,\source{being-glue} and its finding is uncomfortable and correct: the
work is necessary, it is invisible to promotion, and if you only do it you only
get better at doing it.

For an interview this converts into a rule about what to say.

\begin{kittable}{@{}XX@{}}
\textbf{Describes effort} & \textbf{Describes scope} \\
\midrule
``I kept three teams aligned'' & ``I decided the integration contract the three teams built against'' \\
``I wrote a lot of the migration'' & ``I chose expand-and-contract over a cutover, and said what would have changed my mind'' \\
``I mentored several engineers'' & ``Two engineers now review this area without me, and here is what I had to write down for that to be true'' \\
``I fixed a lot of incidents'' & ``I changed the thing that generated the class of incident'' \\
\end{kittable}

The left column is true, sympathetic, and gets assessed at the level below.
\judgement{This is the single highest-leverage editing pass to run over
everything you plan to say: for each sentence, ask whether it describes what
you did or what changed because you did it.}

\section{The archetypes, and the published disagreement}

There is a widely-used taxonomy of Staff roles \dash{} the technical lead, the
architect, the solver, the right hand \dash{} and it is genuinely useful
vocabulary.\source{staff-archetypes}

It is also publicly disputed by practitioners, on grounds worth knowing.
\reported{One argument is that targeting an archetype is bad practical advice,
because influence accumulates from being useful over time and cannot be aimed
at directly, and that the defining trait is accepting accountability for
outcomes you do not fully control.}{archetypes-rebuttal-goedecke} \reported{A
second argues the archetypes function as anti-patterns when treated as
goals.}{archetypes-rebuttal-ewerlof}

\begin{kitmethod}
The usable position, which both sides support:

\textbf{Use the archetypes as vocabulary for describing what you already did},
not as a target to grow into. Saying ``that piece of work was solver-shaped:
nobody owned the problem, it had been open for two quarters, and I was given no
authority over the teams whose code it was'' is a precise, checkable sentence.
Saying ``I am an architect archetype'' is a claim about identity that an
interviewer cannot verify and did not ask for.
\end{kitmethod}

\section{Influence without authority is the assessed skill}

At this level you will be asked, in some form, about a time you needed somebody
who did not report to you, and had no lever, and needed them to change
direction anyway.

The answers that land have a shape. They name what the other person was
optimising for, which is almost always reasonable from where they stood; they
name what you changed about the problem rather than about the person \dash{} a
measurement produced, a cost made visible, a smaller first step proposed; and
they say honestly whether it worked. \judgement{An answer in which you were
right, they were wrong, and they eventually came round is the weakest version
of this story, because it describes persistence rather than influence.}

\section{What this means for a transition candidate}

You are asking to be assessed at Staff in a domain where you have no title.
That combination has one honest consequence and it is better to say it than to
have it found.

Your Staff evidence comes from the work you have actually done, whatever domain
it was in, because scope, ambiguity and influence do not belong to a
technology. Your AI-engineering evidence comes from artefacts, because you have
no job history to point at. These are two separate claims resting on two
separate bodies of evidence, and Part~\ref{part:evidence} is about assembling
the second without weakening the first.

\begin{drill}
Take the three pieces of work you are most likely to be asked about. For each,
write one sentence answering ``who decided this was the problem?''. If the
answer is somebody else for all three, you have found what to work on, and it
is not an interview problem.
\end{drill}

\section*{What this chapter does not claim}

It does not claim the archetype taxonomy is wrong. It claims the disagreement
about how to use it is real and published, and that presenting one side as
settled would be the kind of confident summary this book is trying to avoid.

It does not offer a figure for how many engineers reach Staff, or what the level
is worth relative to Senior. Both figures circulate widely and both trace only
to recruiting content; they are named and refused in Appendix~\ref{app:refused}.
